\chapter{関連研究}
\label{sec:related_work}

\section{文献を整理する観点}
関連研究の章では，既存研究の内容を列挙するだけでなく，自分の研究でどのような判断を行うかに
結び付けて説明する．例えば，検索手法，評価資料，分析方法という観点に分ければ，
採用する基準システムや評価条件を選んだ理由を述べやすくなる．研究課題ごとに整理する場合には，
各課題について何が分かっており，何が未検討であるかを示すとよい．

文献を紹介する際には，対象とする問題，使用した証拠，適用条件を区別する．
同じ手法名が使われていても，対象となる文書や検索要求が異なれば，結果の意味は変わり得る．
そのため，数値だけを抜き出して優劣を判断するのではなく，測定した条件も合わせて記載する必要がある．
これは，後で自分の実験との共通点や相違点を説明するためにも役立つ．

\section{引用の例}
研究室のウェブサイトの項目\cite{RSL}は，参考文献にURLを表示する例として使用している．
日本語論文の項目\cite{exampleArticle}と英語の会議論文の項目\cite{exampleConference}は，
参考文献の書式を確認するための架空の文献である．これらは本文の方法論的な説明を裏付ける資料ではない．
実際の執筆時には，内容を確認した実在の文献に置き換え，引用した文と出典の対応を確かめること．

引用は，その文献がどの説明を支えているかが分かる位置に置く．複数の研究をまとめて述べる場合でも，
すべての文献が同じ条件や結論を扱っているとは限らない．重要な違いがある場合には，
それを文章で説明する．また，書誌情報の正確さと，本文の説明が原文に忠実であることは別の確認事項であり，
一方が正しくても他方が保証されるわけではない．

% Demonstration-only break.
\newpage
\section{比較可能性の確認}
異なる研究の得点を比較する前に，文書集合，要求の選定方法，判定の範囲，評価指標を確認する．
例えば，短い事実確認型の要求を集めた資料と，広い情報収集を目的とする要求を集めた資料では，
上位の結果に求められる内容が異なる可能性がある．得点の差が手法だけによって生じたと考えることはできない．

比較表は設定の違いを簡潔に示すのに適しているが，表だけで解釈を終えてはいけない．
その違いが自分の研究課題にどのように関係するかを説明する必要がある．
従来研究の制約を指摘する場合にも，その研究が本来何を明らかにするために設計されたかを考慮する．
異なる目的に沿った設計を，単純に欠点として扱わないことが重要である．

\section{本研究との関係}
本例では，追加の順位付け処理による変化を，検索要求ごとに対応を付けて調べる．
関連研究の整理は，この手順を選ぶ理由と，検討すべき代替案を明確にする役割を持つ．
基準システムは，単に実装しやすいという理由だけでなく，比較の目的に合っているかという観点から選ぶ．
設定や必要なデータを再現できることも，実験を解釈する上で重要な条件である．

次の章では，これらの観点を具体的な比較条件に落とし込む．
使用する入力をそろえ，変更する部分を限定し，評価結果を集計する方法を先に定める．
関連研究と実験の説明をこのようにつなぐことで，文献調査が実際の研究設計にどのような影響を与えたかを示せる．
なお，以下の条件や数値は，いずれも本テンプレートの表示確認のための例である．
